% LTeX: language=de

\subsubsection{Aufgaben}

\begin{Exercise}{S. 107/3}{}
\begin{enumerate}
    \item [3.]\begin{enumerate}[label=\alph*),leftmargin=*]
        \item[g)] {
            $f(x) = \arctan\left(\dfrac{1}{x+2}\right)$; \qquad $u(x) = \dfrac{1}{x+2}$; \qquad $z(x) = 1$; \qquad $n(x) = x+2$ \\[-2.0ex]
            \begin{itemize}[leftmargin=*]
                \item Definitionsmenge: \\[1.0ex]{
                    $n(x) \neq 0:$\\[-4.0ex]
                    \begin{align*}
                        n(x) &= 0 &&\\[2.0ex]
                        x + 2 &= 0 &&| \quad -2\\[1.5ex]
                        x_{1} &= -2 \ \text{\textit{(e)}}&&
                    \end{align*}
                    $\Rightarrow \quad \uuline{D_{f} = \mathbb{R} \setminus \left\{ -2 \right\}}$;
                    }
                \item Symmetrie: \\[1.0ex]{
                    $D_{f}$ nicht symmetrisch \quad $\Rightarrow$ \quad \uuline{keine Symmetrie vorhanden}\\[-2.0ex]
                    }
                \item Nullstellen: \\[1.0ex]{
                    $f(x) = 0$ \qquad $\Rightarrow$ \qquad $u(x) = 0$\\[-2.0ex]
                    \begin{align*}
                        z(x) > 0 \quad \Rightarrow \quad \text{keine Nullstellen}
                    \end{align*}
                    }
                \item Grenzverhalten: \\[-3.0ex]{
                    \begin{align*}
                        \lim\limits_{x \to -\infty} \arctan\left(\overbrace{\left[\dfrac{1}{\underbrace{x+2}_{\textcolor{myg}{-\infty}}}\right]}^{\textcolor{myg}{0^{-}}}\right) &= 0^{-}; \qquad \qquad \lim\limits_{x \to +\infty} \arctan\left(\overbrace{\left[\dfrac{1}{\underbrace{x+2}_{\textcolor{myg}{+\infty}}}\right]}^{\textcolor{myg}{0^{+}}}\right) &= 0^{+} &&\\[2.0ex]
                        \lim\limits_{x \to -2^{-}} \arctan\left(\overbrace{\left[\dfrac{1}{\underbrace{x+2}_{\textcolor{myg}{0^{-}}}}\right]}^{\textcolor{myg}{-\infty}}\right) &= - \dfrac{\pi}{2}^{+}; \qquad \qquad \lim\limits_{x \to -2^{+}} \arctan\left(\overbrace{\left[\dfrac{1}{\underbrace{x+2}_{\textcolor{myg}{0^{+}}}}\right]}^{\textcolor{myg}{+\infty}}\right) &= \dfrac{\pi}{2}^{-} && \\[-3.0ex]
                    \end{align*}
                    }
                \item Asymptoten: \\[-2.0ex]{
                    \begin{itemize}
                        \item Waagerechte Asymptoten: \quad {
                            \vspace{-5.0ex}
                            \begin{align*}
                                h_{1}(x): \quad &y = 0
                            \end{align*}
                            }
                    \end{itemize}
                    }
                \item Monotonie: \\[-3.0ex]{
                    \begin{align*}
                        f'(x) &= \dfrac{1}{1 + \left(\dfrac{1}{x+2}\right)^{2}} \cdot \dfrac{\left(x+2\right) \cdot 0 - 1 \cdot 1}{\left(x+2\right)^{2}} &&\\[1.5ex]
                        &= \dfrac{1}{1 + \dfrac{1}{\left(x+2\right)^{2}}} \cdot \dfrac{-1}{\left(x+2\right)^{2}} &&\\[1.5ex]
                        &= \dfrac{-1}{1 + \left(x+2\right)^{2}} &&\\[-2.0ex]
                    \end{align*}
                    $z(x) = -1$; \qquad $n(x) = 1 + \left(x+2\right)^{2}$\\[1.0ex]
                    $f'(x) = 0 \qquad \Rightarrow \qquad z(x) < 0 \qquad \Rightarrow$ \quad keine Nullstellen\\[1.0ex]
                    \begin{minipage}[t]{0.4\linewidth}
                        \begin{align*}
                            n(x) &= 0 \\[2.0ex]
                            1 + \left(x+2\right)^{2} &= 0 &&| \quad -1\\[1.5ex]
                            \left(x+2\right)^{2} &= -1 &&| \quad \pm \sqrt[2]{\phantom{x}}\\[1.5ex]
                            & \textcolor{myr}{\text{n. d. in } \mathbb{R}}\\[-1.5ex]
                        \end{align*}
                    \end{minipage}
                    \hfill
                    \begin{minipage}[t]{0.55\linewidth}
                        \raisebox{-\height}{
                        \begin{tabularx}{0.95\linewidth}{c|X|X|X}
                            & $-\infty < x < -2$ & $x = -2$ & $-2 < x < \infty$ \\ \hline
                            $z(x)$ & $-$ & $-$ & $-$ \\ \hline
                            $n(x)$ & $+$ & $+$ & $+$ \\ \hline
                            $f'(x)$ & $-$ & $-$ & $-$ \\ \hline
                            $G_{f}$ & $\searrow$ &  & $\searrow$ \\
                        \end{tabularx}
                        }  
                    \end{minipage}
                    $\Rightarrow \qquad G_{f}$ ist smf in $\left] \ -\infty; -2 \ \right[$ \quad sowie \quad $\left] \ -2; \infty \ \right[$\\[-2.0ex]
                    }
                \item Krümmung: \\[-2.0ex]{
                    \begin{align*}
                        f''(x) &= \dfrac{\left[1 + \left(x+2\right)^{2}\right] \cdot 0 - \left(-1\right) \cdot 2 \cdot \left(x+2\right) \cdot 1}{\left[1 + \left(x+2\right)^{2}\right]^{2}} &&\\[1.5ex]
                        &= \dfrac{2 \cdot \left(x + 2\right)}{1 + \left(x+2\right)^{4}} &&\\[-1.5ex]
                    \end{align*}
                    $z(x) = 2 \cdot \left(x + 2\right)$; \qquad $n(x) = 1 + \left(x+2\right)^{4}$\\[1.0ex]
                    $f''(x) = 0 \qquad \Rightarrow \qquad z(x) > 0 \qquad \Rightarrow$ \quad keine Nullstellen\\[1.0ex]
                    \begin{minipage}[t]{0.4\linewidth}
                        \vspace{-2.0ex}
                        \begin{align*}
                            n(x) &= 0 \\[2.0ex]
                            1 + \left(x+2\right)^{4} &= 0 &&| \quad - 1\\[1.5ex]
                            \left(x+2\right)^{4} &= -1 &&| \quad \pm \sqrt[4]{\phantom{x}}\\[1.5ex]
                            & \textcolor{myr}{\text{n. d. in } \mathbb{R}}\\[-1.5ex]
                        \end{align*}
                    \end{minipage}
                    \hfill
                    \begin{minipage}[t]{0.55\linewidth}
                        \vspace{-2.0ex}
                        \begin{align*}
                            z(x) &= 0\\[2.0ex]
                            2 \cdot \left(x + 2\right) &= 0 &&| \quad :2\\[1.5ex]
                            x + 2 &= 0 &&| \quad -2\\[1.5ex]
                            x_{1} &= -2 \ \text{\textit{(e)}}&&\\[-2.0ex]
                        \end{align*}
                    \end{minipage}
                    \raisebox{-\height}{
                    \begin{tabularx}{0.95\linewidth}{c|X|X|X}
                        & $-\infty < x < -2$ & $x = -2$ & $-2 < x < \infty$ \\ \hline
                        $z(x)$ & $-$ & $0$ & $+$ \\ \hline
                        $n(x)$ & $+$ & $+$ & $+$ \\ \hline
                        $f''(x)$ & $-$ & $0$ & $+$ \\ \hline
                        $G_{f}$ & r. g. &  & l. g. \\
                    \end{tabularx}
                    }\\[2.0ex]
                    $\Rightarrow \qquad G_{f}$ ist rechtsgekrümmt in $\left] \ -\infty; -2 \ \right[$ \\[1.0ex]
                    $\Rightarrow \qquad G_{f}$ ist linksgekrümmt in $\left] \ -2; +\infty \ \right[$\\[-2.0ex]
                    }
                \item Extrem- und Wendepunkte: \qquad keine vorhanden
                \item Zeichnung: \\[-2.0ex]{
                    \raisebox{-\height}{
                    \begin{tikzpicture}[
                            declare function={
                                f_f(\x) = rad(atan(1/((\x*\xScale) + 2))))*(1/\yScale);
                            }
                        ]

                        % Set variables
                        \pgfmathsetmacro{\axisPosWidth}{4.0};
                        \pgfmathsetmacro{\axisNegWidth}{-8.0};
                        \pgfmathsetmacro{\axisPosHeight}{2.0};
                        \pgfmathsetmacro{\axisNegHeight}{-2.0};
                        \pgfmathsetmacro{\xIncrement}{1.0};
                        \pgfmathsetmacro{\yIncrement}{1.0};
                        \pgfmathsetmacro{\xScale}{1.0};
                        \pgfmathsetmacro{\yScale}{1.0};
                        \pgfmathsetmacro{\gridxIncrement}{0.5};
                        \pgfmathsetmacro{\gridyIncrement}{0.5};

                        % Axis/grid limits
                        \pgfmathsetmacro{\xUpperLimitAxis}{\xIncrement * floor(\axisPosWidth / \xIncrement)}
                        \pgfmathsetmacro{\xLowerLimitAxis}{\xIncrement * ceil(\axisNegWidth / \xIncrement)}
                        \pgfmathsetmacro{\yUpperLimitAxis}{\yIncrement * floor(\axisPosHeight /\yIncrement)}
                        \pgfmathsetmacro{\yLowerLimitAxis}{\yIncrement * ceil(\axisNegHeight / \yIncrement)}

                        \pgfmathsetmacro{\xUpperLimitGrid}{\gridxIncrement * floor(\axisPosWidth / \gridxIncrement)}
                        \pgfmathsetmacro{\xLowerLimitGrid}{\gridxIncrement * ceil(\axisNegWidth / \gridxIncrement)}
                        \pgfmathsetmacro{\yUpperLimitGrid}{\gridyIncrement * floor(\axisPosHeight / \gridyIncrement)}
                        \pgfmathsetmacro{\yLowerLimitGrid}{\gridyIncrement * ceil(\axisNegHeight / \gridyIncrement)}

                        % Axis environment (PGFPlots)
                        \begin{axis}[
                            xmin=\xLowerLimitGrid, xmax=\xUpperLimitGrid,
                            ymin=\yLowerLimitGrid, ymax=\yUpperLimitGrid,
                            axis lines=middle,
                            axis line style={->, >=stealth},
                            grid=both,
                            xlabel={$x$},
                            ylabel={$y$},
                            xlabel style={anchor=west, at={(current axis.right of origin)}},
                            ylabel style={anchor=south, at={(current axis.above origin)}},
                            xtick={\xLowerLimitAxis,\xLowerLimitAxis+1,...,\xUpperLimitAxis},
                            ytick={\yLowerLimitAxis,\yLowerLimitAxis+1,...,\yUpperLimitAxis},
                            xticklabel={\pgfmathparse{\tick*\xScale}\pgfmathprintnumber{\pgfmathresult}},
                            yticklabel={\pgfmathparse{\tick*\yScale}\pgfmathprintnumber{\pgfmathresult}},
                            tick label style={font=\scriptsize},
                            x=1cm,
                            y=1cm,
                            minor tick num=1,
                            scale only axis,
                            grid=both,
                            restrict y to domain=\yLowerLimitGrid:\yUpperLimitGrid,
                        ]
                            
                        % Graph limits
                        \pgfmathsetmacro{\xLowerLimitGraph}{\xLowerLimitAxis}
                        \pgfmathsetmacro{\xUpperLimitGraph}{\xUpperLimitAxis}
                        
                        % G_f
                        \pgfmathsetmacro{\Gfax}{0.5*(1/\xScale)}
                        \pgfmathsetmacro{\Gfay}{f_f(\Gfax)}
                        \addplot[thick, myb, samples=1000, domain=\xLowerLimitGraph:-2.02*(1/\xScale)]
                            {f_f(x)};
                        \addplot[thick, myb, samples=1000, domain=-1.98*(1/\xScale):\xUpperLimitGraph]
                            {f_f(x)};
                        \node[anchor=south west] at (axis cs:{\Gfax},{\Gfay}) {\textcolor{myb}{$G_{f}$}};
    
                        % Definitionslücken
                        \pgfmathsetmacro{\Dfax}{-1.98*(1/\xScale)}
                        \pgfmathsetmacro{\Dfay}{f_f(\Dfax)}
                        \addplot[only marks, mark=o, mark size=2pt, thick] coordinates {(\Dfax,\Dfay)};
    
                        \pgfmathsetmacro{\Dfbx}{-2.02*(1/\xScale)}
                        \pgfmathsetmacro{\Dfby}{f_f(\Dfbx)}
                        \addplot[only marks, mark=o, mark size=2pt, thick] coordinates {(\Dfbx,\Dfby)};
    
                        % h_1 für y1
                        \pgfmathsetmacro{\Gtby}{0*(1/\yScale)}
                        \addplot[thick, myorange] coordinates {(\xLowerLimitGraph,\Gtby) (\xUpperLimitGraph,\Gtby)};
                        \node[anchor=north west] at (axis cs:{0},{\Gtby}) {\textcolor{myorange}{$G_{h_{1}}$}};

                        \end{axis}
                    \end{tikzpicture}
                    }
                    }
            \end{itemize}
            }
        \item[h)] {
            $f(x) = \arctan\left(\sqrt[2]{\dfrac{1}{x^{2}+4}}\right)$; \\[1.5ex]
            $u(x) = \sqrt[2]{\dfrac{1}{x^{2}+4}}$; \qquad $v(x) = \dfrac{1}{x^{2}+4}$; \qquad $z(x) = 1$; \qquad $n(x) = x^{2}+4$ \\[-2.0ex]
            \begin{itemize}[leftmargin=*]
                \item Definitionsmenge: \\[1.0ex]{
                    \begin{minipage}[t]{0.5\linewidth}
                        $n(x) \neq 0:$\\[-4.0ex]
                        \begin{align*}
                            n(x) &= 0 &&\\[2.0ex]
                            x^{2} + 4 &= 0 &&| \quad -4\\[1.5ex]
                            x^{2} &= -4 &&| \quad \pm \sqrt[2]{\phantom{x}}\\[1.5ex]
                            & \textcolor{myr}{\text{n. d. in } \mathbb{R}}\\[-1.5ex]
                        \end{align*}
                    \end{minipage}
                    \begin{minipage}[t]{0.5\linewidth}
                        $v(x) \geq 0:$\\[-4.0ex]
                        \begin{align*}
                            v(x) &= 0 \qquad \Rightarrow \quad z(x) > 0&&
                        \end{align*}
                        $\Rightarrow$ \quad keine Nullstellen\\[-1.0ex]
                        \raisebox{-\height}{
                        \begin{tabularx}{0.95\linewidth}{c|X|X|X}
                            & $-\infty < x < -4$ & $x = -4$ & $-4 < x < \infty$ \\ \hline
                            $z(x)$ & $+$ & $+$ & $+$ \\ \hline
                            $n(x)$ & $+$ & $+$ & $+$ \\ \hline
                            $v(x)$ & $+$ & $+$ & $+$ \\
                        \end{tabularx}
                        }\\[2.0ex]
                    \end{minipage}
                    $\Rightarrow \quad \uuline{D_{f} = \mathbb{R}}$\\[-2.0ex]
                    }
                \item Symmetrie: \\[-4.0ex]{
                    \begin{align*}
                        f(-x) &= f(x) &&\\[2.0ex]
                        \arctan\left(\sqrt[2]{\dfrac{1}{\left(-x\right)^{2}+4}}\right) &= \arctan\left(\sqrt[2]{\dfrac{1}{x^{2}+4}}\right) &&\\[1.5ex]
                        \arctan\left(\sqrt[2]{\dfrac{1}{x^{2}+4}}\right) &= \arctan\left(\sqrt[2]{\dfrac{1}{x^{2}+4}}\right) \quad \text{\textit{(wahr)}} &&
                    \end{align*}
                    $\Rightarrow$ \quad $G_f$ ist achsensymmetrischzur y-Achse\\[-2.0ex]
                    }
                \item Nullstellen: \qquad $u(x) > 0$ \qquad $\Rightarrow$ \qquad $f(x) > 0$ \qquad $\Rightarrow$ \quad keine Nullstellen\\[-2.0ex]
                \item Grenzverhalten: \\[-3.0ex]{
                    \begin{align*}
                        \lim\limits_{x \to -\infty} \arctan\left(\overbrace{\sqrt[2]{\dfrac{1}{\underbrace{x^{2}+4}_{\textcolor{myg}{+\infty}}}}}^{\textcolor{myg}{0^{+}}}\right) &= 0^{+}; \qquad \qquad \lim\limits_{x \to +\infty} \arctan\left(\overbrace{\sqrt[2]{\dfrac{1}{\underbrace{x^{2}+4}_{\textcolor{myg}{+\infty}}}}}^{\textcolor{myg}{0^{+}}}\right) &= 0^{+} &&\\[-5.0ex]
                    \end{align*}
                    }
                \item Asymptoten: \\[-3.0ex]{
                    \begin{itemize}
                        \item Waagerechte Asymptoten: \quad {
                            \vspace{-5.0ex}
                            \begin{align*}
                                h_{1}(x): \quad &y = 0
                            \end{align*}
                            }
                    \end{itemize}
                    }
                \item Monotonie: \\[-2.0ex]{
                    \begin{align*}
                        f'(x) &= \dfrac{1}{1 + \left(\sqrt[2]{\dfrac{1}{x^{2}+4}}\right)^{2}} \cdot \dfrac{1}{2} \cdot \dfrac{1}{\sqrt[2]{\dfrac{1}{x^{2}+4}}} \cdot \dfrac{\left(x^{2} + 4\right) \cdot 0 - 1 \cdot 2x}{\left(x^{2} + 4\right)^{2}} &&\\[1.5ex]
                        &= \dfrac{1}{1 + \dfrac{1}{x^{2}+4}} \cdot \dfrac{1}{2} \cdot \sqrt[2]{\dfrac{x^{2}+4}{1}} \cdot \dfrac{-2x}{\left(x^{2} + 4\right)^{2}} &&\\[1.5ex]
                        &= \dfrac{1}{1 + \dfrac{1}{x^{2}+4}} \cdot \sqrt[2]{x^{2}+4} \cdot \dfrac{-x}{\left(x^{2} + 4\right)^{2}} &&\\[1.5ex]
                        &= \dfrac{1}{1 + \dfrac{1}{x^{2}+4}} \cdot \dfrac{-x}{\left(x^{2} + 4\right)^{\frac{3}{2}}} &&\\[1.5ex]
                        &= \dfrac{-x}{\left(x^{2} + 4\right)^{\frac{3}{2}} + \left(x^{2} + 4\right)^{\frac{1}{2}}} &&\\[1.5ex]
                        &= \dfrac{-x}{\left(x^{2} + 4\right)^{\frac{1}{2}} \cdot \left[\left(x^{2} + 4\right)^{1} + 1 \right]} &&\\[1.5ex]
                        &= \dfrac{-x}{\left(x^{2} + 4\right)^{\frac{1}{2}} \cdot \left(x^{2} + 5 \right)} &&\\[-2.0ex]
                    \end{align*}
                    $z(x) = -x$; \\[1.0ex]
                    $n(x) = \left(x^{2} + 4\right)^{\frac{1}{2}} \cdot \left(x^{2} + 5 \right)$; \qquad $u(x) = \left(x^{2} + 4\right)^{\frac{1}{2}}$; \qquad $v(x) = x^{2} + 5$\\[1.5ex]
                    $f'(x) = 0:$\\
                    \begin{minipage}[t]{0.20\linewidth}
                        \vspace{-2.0ex}
                        \begin{align*}
                            z(x) &= 0 \\[2.0ex]
                            -x &= 0 &&\\[1.5ex]
                            x_{1} &= 0 \ \text{\textit{(e)}}&&\\[-2.0ex]
                        \end{align*}
                    \end{minipage}
                    \hfill
                    \begin{minipage}[t]{0.35\linewidth}
                        \vspace{-2.0ex}
                        \begin{align*}
                            u(x) &= 0 \\[2.0ex]
                            \left(x^{2} + 4\right)^{\frac{1}{2}} &= 0 &&| \quad ^{2}\\[1.5ex]
                            x^{2} + 4 &= 0 &&| \quad - 4\\[1.5ex]
                            x^{2} &= -4 &&| \quad \pm \sqrt[2]{\phantom{x}}\\[1.5ex]
                            & \textcolor{myr}{\text{n. d. in } \mathbb{R}}\\[-1.5ex]
                        \end{align*}
                    \end{minipage}
                    \hfill
                    \begin{minipage}[t]{0.25\linewidth}
                        \vspace{-2.0ex}
                        \begin{align*}
                            v(x) &= 0 \\[2.0ex]
                            x^{2} + 5 &= 0 &&| \quad -5\\[1.5ex]
                            x^{2} &= -5 &&| \quad \pm \sqrt[2]{\phantom{x}}\\[1.5ex]
                            & \textcolor{myr}{\text{n. d. in } \mathbb{R}}\\[-1.5ex]
                        \end{align*}
                    \end{minipage}
                    \raisebox{-\height}{
                    \begin{tabularx}{0.95\linewidth}{c|X|X|X}
                        & $-\infty < x < 0$ & $x = 0$ & $0 < x < \infty$ \\ \hline
                        $z(x)$ & $+$ & $0$ & $-$ \\ \hline
                        $u(x)$ & $+$ & $+$ & $+$ \\ \hline
                        $v(x)$ & $+$ & $+$ & $+$ \\ \hline
                        $n(x)$ & $+$ & $+$ & $+$ \\ \hline
                        $f'(x)$ & $+$ & $0$ & $-$ \\ \hline
                        $G_{f}$ & $\nearrow$ & HoP & $\searrow$ \\
                    \end{tabularx}
                    }\\[2.0ex]
                    $\Rightarrow \qquad G_{f}$ ist sms in $\left] \ -\infty; 0 \ \right]$ \\[1.0ex]
                    $\Rightarrow \qquad G_{f}$ ist smf in $\left[ \ 0; \infty \ \right[$\\[-2.0ex]
                    }
                \item Krümmung: \\[-2.0ex]{
                    \begin{align*}
                        f''(x) &= \dfrac{\left(x^{2} + 4\right)^{\frac{1}{2}} \cdot \left(x^{2} + 5 \right) \cdot \left(-1\right) - \left(-x\right) \cdot \left[\dfrac{1}{2} \cdot \dfrac{1}{\left(x^{2} + 4\right)^{\frac{1}{2}}} \cdot 2x \cdot \left(x^{2} + 5\right) + \left(x^{2} + 4\right)^{\frac{1}{2}} \cdot 2x\right]}{\left[\left(x^{2} + 4\right)^{\frac{1}{2}} \cdot \left(x^{2} + 5 \right)\right]^{2}} &&\\[2.0ex]
                        &= \dfrac{\left(x^{2} + 4\right)^{\frac{1}{2}} \cdot \left(x^{2} + 5 \right) \cdot \left(-1\right) + x \cdot \left[\dfrac{1}{\left(x^{2} + 4\right)^{\frac{1}{2}}} \cdot x \cdot \left(x^{2} + 5\right) + \left(x^{2} + 4\right)^{\frac{1}{2}} \cdot 2x\right]}{\left[\left(x^{2} + 4\right)^{\frac{1}{2}}\right]^{2} \cdot \left(x^{2} + 5 \right)^{2}} &&\\[2.0ex]
                        &= \dfrac{\left(x^{2} + 4\right) \cdot \left(x^{2} + 5 \right) \cdot \left(-1\right) + x \cdot \left[1 \cdot x \cdot \left(x^{2} + 5\right) + \left(x^{2} + 4\right) \cdot 2x\right]}{\left(x^{2} + 4\right)^{\frac{1}{2}} \cdot \left[\left(x^{2} + 4\right)^{\frac{1}{2}}\right]^{2} \cdot \left(x^{2} + 5 \right)^{2}} &&\\[2.0ex]
                        &= \dfrac{\left(x^{4} + 5 x^{2} + 4x^{2} + 20\right) \cdot \left(-1\right) + x \cdot \left[x^{3} + 5x + 2x^{3} + 8x\right]}{\left(x^{2} + 4\right)^{\frac{3}{2}} \cdot \left(x^{2} + 5 \right)^{2}} &&\\[-2.0ex]
                    \end{align*}
                    \begin{align*}
                        f''(x) &= \dfrac{\left(x^{4} + 9 x^{2} + 20\right) \cdot \left(-1\right) + x \cdot \left[3x^{3} + 13x\right]}{\left(x^{2} + 4\right)^{\frac{3}{2}} \cdot \left(x^{2} + 5 \right)^{2}} &&\\[2.0ex]
                        &= \dfrac{-x^{4} - 9 x^{2} - 20 + 3x^{4} + 13x^{2}}{\left(x^{2} + 4\right)^{\frac{3}{2}} \cdot \left(x^{2} + 5 \right)^{2}} &&\\[2.0ex]
                        &= \dfrac{2x^{4} + 4x^{2} - 20}{\left(x^{2} + 4\right)^{\frac{3}{2}} \cdot \left(x^{2} + 5 \right)^{2}} &&\\[2.0ex]
                        &= \dfrac{2 \cdot \left(x^{4} + 2x^{2} - 10\right)}{\left(x^{2} + 4\right)^{\frac{3}{2}} \cdot \left(x^{2} + 5 \right)^{2}} &&\\[-2.0ex]
                    \end{align*}
                    $z(x) = 2 \cdot \left(x^{4} + 2x^{2} - 10\right)$; \\[1.0ex]
                    $n(x) = \left(x^{2} + 4\right)^{\frac{3}{2}} \cdot \left(x^{2} + 5 \right)^{2}$; \qquad $u(x) = \left(x^{2} + 4\right)^{\frac{3}{2}}$; \qquad $v(x) = \left(x^{2} + 5 \right)^{2}$\\[1.0ex]
                    $f''(x) = 0:$\\[-1.0ex]
                    \vspace{-2.0ex}
                    \begin{align*}
                        z(x) &= 0 \\[2.0ex]
                        2 \cdot \left(x^{4} + 2x^{2} - 10\right) &= 0 &&| \quad :2\\[1.5ex]
                        x^{4} + 2x^{2} - 10 &= 0; && z = x^{2}\\[-1.5ex]
                    \end{align*}
                    \begin{align*}
                        z^{2} + 2z - 10 &= 0 &&\\[2.0ex]
                        z_{1,2} &= \dfrac{-2 \pm \sqrt[2]{2^{2} - 4 \cdot 1 \cdot \left(-10\right)}}{2 \cdot 1} && \\[2.0ex]
                        &= \dfrac{-2 \pm 2\sqrt[2]{11}}{2} && \\[-2.0ex]
                    \end{align*}
                    \begin{minipage}[t]{0.35\linewidth}
                        \vspace{-2.0ex}
                        \begin{align*}
                            z_{1} &= -1 - \sqrt[2]{11} &&\\[2.0ex]
                            x^{2} &= -1 - \sqrt[2]{11} &&| \quad \pm \sqrt[2]{\phantom{x}}\\[1.5ex]
                            & \textcolor{myr}{\text{n. d. in } \mathbb{R}} && \\[-2.0ex]
                        \end{align*}
                    \end{minipage}
                    \hfill
                    \vrule
                    \hfill
                    \begin{minipage}[t]{0.60\linewidth}
                        \vspace{-2.0ex}
                        \begin{align*}
                            z_{2} &= -1 + \sqrt[2]{11} &&\\[2.0ex]
                            x^{2} &= -1 + \sqrt[2]{11} &&| \quad \pm \sqrt[2]{\phantom{x}} && \\[-1.5ex]
                        \end{align*}
                        $x_{1} = -\sqrt{-1 + \sqrt{11}} \ \text{\textit{(e)}}; \quad x_{1} \approx \SI{-1.52}{}$\\[2.0ex]
                        $x_{2} = \sqrt{-1 + \sqrt{11}} \ \text{\textit{(e)}}; \quad x_{2} \approx \SI{1.52}{}$\\[-1.0ex]
                    \end{minipage}\\[2.0ex]
                    $n(x) = 0:$\\[-1.0ex]
                    \begin{minipage}[t]{0.45\linewidth}
                        \vspace{-2.0ex}
                        \begin{align*}
                            u(x) &= 0 \\[2.0ex]
                            \left(x^{2} + 4\right)^{\frac{3}{2}} &= 0 &&| \quad ^{\frac{2}{3}}\\[1.5ex]
                            x^{2} + 4 &= 0 &&| \quad -4\\[1.5ex]
                            x^{2} &= -4 &&| \quad \pm \sqrt[2]{\phantom{x}}\\[1.5ex]
                            & \textcolor{myr}{\text{n. d. in } \mathbb{R}}\\[-1.5ex]
                        \end{align*}
                    \end{minipage}
                    \hfill
                    \begin{minipage}[t]{0.45\linewidth}
                        \vspace{-2.0ex}
                        \begin{align*}
                            v(x) &= 0 \\[2.0ex]
                            \left(x^{2} + 5 \right)^{2} &= 0 &&| \quad \sqrt[2]{\phantom{x}}\\[1.5ex]
                            x^{2} + 5 &= 0 &&| \quad -5\\[1.5ex]
                            x^{2} &= -5 &&| \quad \pm \sqrt[2]{\phantom{x}}\\[1.5ex]
                            & \textcolor{myr}{\text{n. d. in } \mathbb{R}}\\[-1.5ex]
                        \end{align*}
                    \end{minipage}
                    \raisebox{-\height}{
                    \begin{tabularx}{0.95\linewidth}{c|X|X|X|X|X}
                        & $-\infty < x < x_{1}$ & $x = x_{1}$ & $x_{1} < x < x_{2}$ & $x = x_{2}$ & $x_{2} < x < \infty$ \\ \hline
                        $z(x)$ & $+$ & $0$ & $-$ & $0$ & $+$ \\ \hline
                        $u(x)$ & $+$ & $+$ & $+$ & $+$ & $+$ \\ \hline
                        $v(x)$ & $+$ & $+$ & $+$ & $+$ & $+$ \\ \hline
                        $n(x)$ & $+$ & $+$ & $+$ & $+$ & $+$ \\ \hline
                        $f''(x)$ & $-$ & $0$ & $+$ & $0$ & $-$ \\ \hline
                        $G_{f}$ & l. g. & $W_{1}$ & r. g. & $W_{2}$ & l. g. \\
                    \end{tabularx}
                    }\\[2.0ex]
                    $\Rightarrow \qquad G_{f}$ ist linksgekrümmt in $\left] \ -\infty; -\sqrt{-1 + \sqrt{11}} \ \right]$ sowie $\left[ \ \sqrt{-1 + \sqrt{11}}; +\infty \ \right[$ \\[1.0ex]
                    $\Rightarrow \qquad G_{f}$ ist rechtsgekrümmt in $\left[ \ -\sqrt{-1 + \sqrt{11}}; \sqrt{-1 + \sqrt{11}} \ \right]$\\[-2.0ex]
                    }
                \item Extrem- und Wendepunkte: \\[-3.0ex]{
                    \begin{itemize}
                        \item Hochpunkt: \quad $H_{1}(0 \ / \ \SI{0.46}{})$
                        \item Wendepunkte: \quad $W_{1}(-\sqrt{-1 + \sqrt{11}} \ / \ \SI{0.38}{})$; \quad $W_{2}(\sqrt{-1 + \sqrt{11}} \ / \ \SI{0.38}{})$
                    \end{itemize}
                }
                \item Zeichnung: \\[-2.0ex]{
                    \raisebox{-\height}{
                    \begin{tikzpicture}[
                            declare function={
                                f_f(\x) = rad(atan(sqrt(1/((\x*\xScale)^2 + 4)))))*(1/\yScale);
                            }
                        ]

                        % Set variables
                        \pgfmathsetmacro{\axisPosWidth}{6.0};
                        \pgfmathsetmacro{\axisNegWidth}{-6.0};
                        \pgfmathsetmacro{\axisPosHeight}{3.0};
                        \pgfmathsetmacro{\axisNegHeight}{-1.0};
                        \pgfmathsetmacro{\xIncrement}{1.0};
                        \pgfmathsetmacro{\yIncrement}{1.0};
                        \pgfmathsetmacro{\xScale}{1.0};
                        \pgfmathsetmacro{\yScale}{0.25};
                        \pgfmathsetmacro{\gridxIncrement}{0.5};
                        \pgfmathsetmacro{\gridyIncrement}{0.5};

                        % Axis/grid limits
                        \pgfmathsetmacro{\xUpperLimitAxis}{\xIncrement * floor(\axisPosWidth / \xIncrement)}
                        \pgfmathsetmacro{\xLowerLimitAxis}{\xIncrement * ceil(\axisNegWidth / \xIncrement)}
                        \pgfmathsetmacro{\yUpperLimitAxis}{\yIncrement * floor(\axisPosHeight /\yIncrement)}
                        \pgfmathsetmacro{\yLowerLimitAxis}{\yIncrement * ceil(\axisNegHeight / \yIncrement)}

                        \pgfmathsetmacro{\xUpperLimitGrid}{\gridxIncrement * floor(\axisPosWidth / \gridxIncrement)}
                        \pgfmathsetmacro{\xLowerLimitGrid}{\gridxIncrement * ceil(\axisNegWidth / \gridxIncrement)}
                        \pgfmathsetmacro{\yUpperLimitGrid}{\gridyIncrement * floor(\axisPosHeight / \gridyIncrement)}
                        \pgfmathsetmacro{\yLowerLimitGrid}{\gridyIncrement * ceil(\axisNegHeight / \gridyIncrement)}

                        % Axis environment (PGFPlots)
                        \begin{axis}[
                            xmin=\xLowerLimitGrid, xmax=\xUpperLimitGrid,
                            ymin=\yLowerLimitGrid, ymax=\yUpperLimitGrid,
                            axis lines=middle,
                            axis line style={->, >=stealth},
                            grid=both,
                            xlabel={$x$},
                            ylabel={$y$},
                            xlabel style={anchor=west, at={(current axis.right of origin)}},
                            ylabel style={anchor=south, at={(current axis.above origin)}},
                            xtick={\xLowerLimitAxis,\xLowerLimitAxis+1,...,\xUpperLimitAxis},
                            ytick={\yLowerLimitAxis,\yLowerLimitAxis+1,...,\yUpperLimitAxis},
                            xticklabel={\pgfmathparse{\tick*\xScale}\pgfmathprintnumber{\pgfmathresult}},
                            yticklabel={\pgfmathparse{\tick*\yScale}\pgfmathprintnumber{\pgfmathresult}},
                            tick label style={font=\scriptsize},
                            x=1cm,
                            y=1cm,
                            minor tick num=1,
                            scale only axis,
                            grid=both,
                            restrict y to domain=\yLowerLimitGrid:\yUpperLimitGrid,
                        ]
                            
                        % Graph limits
                        \pgfmathsetmacro{\xLowerLimitGraph}{\xLowerLimitAxis}
                        \pgfmathsetmacro{\xUpperLimitGraph}{\xUpperLimitAxis}
                        
                        % G_f
                        \pgfmathsetmacro{\Gfax}{0.5*(1/\xScale)}
                        \pgfmathsetmacro{\Gfay}{f_f(\Gfax)}
                        \addplot[thick, myb, samples=1000, domain=\xLowerLimitGraph:\xUpperLimitGraph]
                            {f_f(x)};
                        \node[anchor=south west] at (axis cs:{\Gfax},{\Gfay}) {\textcolor{myb}{$G_{f}$}};
    
                        % h_1 für y1
                        \pgfmathsetmacro{\Gtay}{0*(1/\yScale)}
                        \addplot[thick, myorange] coordinates {(\xLowerLimitGraph,\Gtay) (\xUpperLimitGraph,\Gtay)};
                        \node[anchor=north west] at (axis cs:{0},{\Gtay}) {\textcolor{myorange}{$G_{h_{1}}$}};

                        % H_1
                        \pgfmathsetmacro{\Hax}{0}
                        \pgfmathsetmacro{\Hay}{f_f(\Hax)}
                        \pgfmathsetmacro{\HaxLabel}{\Hax*\xScale}
                        \pgfmathsetmacro{\HayLabel}{\Hay*\yScale}
                        \addplot[only marks, mark=x, mark size=3pt] coordinates {(\Hax,\Hay)};
                        \node[anchor=south east, xshift=0ex, yshift=0.25cm] at (axis cs:{\Hax},{\Hay}) {$H_{1}(\num[output-decimal-marker = {.},round-mode=places, round-precision=2]{\HaxLabel}{},\ \num[output-decimal-marker = {.},round-mode=places,round-precision=2]{\HayLabel}{})$};

                        % W_1
                        \pgfmathsetmacro{\Wax}{-sqrt(-1 + sqrt(11))}
                        \pgfmathsetmacro{\Way}{f_f(\Wax)}
                        \pgfmathsetmacro{\WaxLabel}{\Wax*\xScale}
                        \pgfmathsetmacro{\WayLabel}{\Way*\yScale}
                        \addplot[only marks, mark=x, mark size=3pt] coordinates {(\Wax,\Way)};
                        \node[anchor=south east, xshift=0ex, yshift=0ex] at (axis cs:{\Wax},{\Way}) {$W_{1}(\num[output-decimal-marker = {.},round-mode=places, round-precision=2]{\WaxLabel}{},\ \num[output-decimal-marker = {.},round-mode=places,round-precision=2]{\WayLabel}{})$};
    
                        % W_2
                        \pgfmathsetmacro{\Wbx}{sqrt(-1 + sqrt(11))}
                        \pgfmathsetmacro{\Wby}{f_f(\Wax)}
                        \pgfmathsetmacro{\WbxLabel}{\Wbx*\xScale}
                        \pgfmathsetmacro{\WbyLabel}{\Wby*\yScale}
                        \addplot[only marks, mark=x, mark size=3pt] coordinates {(\Wbx,\Wby)};
                        \node[anchor=south west, xshift=0ex, yshift=0ex] at (axis cs:{\Wbx},{\Wby}) {$W_{2}(\num[output-decimal-marker = {.},round-mode=places, round-precision=2]{\WbxLabel}{},\ \num[output-decimal-marker = {.},round-mode=places,round-precision=2]{\WbyLabel}{})$};
                        \end{axis}
                    \end{tikzpicture}
                    }
                    }
            \end{itemize}
            }
    \end{enumerate}
\end{enumerate}
\end{Exercise}